\color{unimain}{\subsection{Максимальная токовая защита стороны НН2(CН)}}
\color{black}

\begin{enumerate}

\item
Направленная максимальная токовая защита на стороне НН2(СН) с комбинированным пуском по напряжению предназначена для резервирования основных защит трансформатора, резервирования отключения КЗ на шинах НН2(СН) и на элементах присоединенных к шинам НН2(СН).

\item
Алгоритм работы МТЗ НН2(СН) приведен на рисунке \ref{pic:MTZ_NN2}. Параметры для настройки МТЗ НН2(СН) приведены в таблице \ref{app_set:mtznn2}.

\medskip
\begin{figure}[H]
\centering
\includegraphics[width=1\textwidth,height=1\textheight,keepaspectratio]{\fbpath/06.01 ДЗТ/МТЗ НН2/img/MTZ_NN2.pdf}
\captionof{figure}{Функциональная схема логики МТЗ НН2(СН)}\label{pic:MTZ_NN2}
\end{figure}

\item
Каждая из ступеней имеет регулируемые значения по току и времени срабатывания. Значения срабатывания по току определяются параметрами \textbf{«Iср МТЗ-1 НН2(СН)»} и \textbf{«Iср МТЗ-2 НН2(СН)»} для
первой и второй ступеней соответственно. Задержки на срабатывание защиты задаются параметрами \textbf{«Тср МТЗ-1 НН2(СН)»} и \textbf{«Тср МТЗ-2 НН2(СН)»}.

\item
Каждая из ступеней защиты может выполнена с пуском по напряжению. Ввод контроля пуска по напряжению в логику работы ступеней осуществляется программными накладками \textbf{«Режим КПОН МТЗ-1 НН2(СН)»} и \textbf{«Режим КПОН МТЗ-2 НН2(СН)»} для первой и второй ступеней соответственно.

\item
Формирование сигнала срабатывания КПОН низкой и средней стороны трансформатора производится с контролем положения выключателей соответствующих сторон.

\item
Каждая из ступеней защиты может выполнена с блокировкой по второй гармонике (отношению уровня второй гармоники к первой). Данная блокировка реализована на базе модуля «ОВ БНТ НН2(СН)». Ввод контроля броска намагничивающего тока программными накладками \textbf{«Режим БНТ МТЗ-1 НН2(СН)»}, \textbf{«Режим БНТ МТЗ-2 НН2(СН)»}.

\item
Для обеспечения несрабатывания МТЗ НН2(СН) при броске тока намагничивания используется перекрестная блокировка по 2 гармонике в фазных токах НН2(СН). Вводится программной накладкой \textbf{«БНТ МТЗ НН2(СН)»}. Блокировка основана на контроле отношения действующего значения второй гармонике к первой токов НН2(СН). Данная блокировка вводиться в работу для каждой ступени независимо с помощью уставок \textbf{«Режим БНТ МТЗ-1 НН2(СН)»} и \textbf{«Режим БНТ МТЗ-2 М»}. Имеется возможность регулировать порог срабатывания блокировки, с помощью уставки \textbf{«I2г/I1г МТЗ НН2(СН)»}. Блокировка МТЗ НН2(СН) по 2 гармонике выводится из действия при величине соответствующего фазного тока ниже 0,04Iном. Также блокировка МТЗ НН2(СН) по 2 гармонике выводится из действия при достижении величины действующего значения тока выше уставки \textbf{«Iср БНТ МТЗ НН2(СН)»}.

\item
Орган направления мощности (рисунок \ref{pic:ONM}) выполнен по 90-градусной схеме сочетания токов и напряжений Ia и Ubc, Ic и Uab. Направление мощности определяется по величине фазового угла между
током и напряжением отдельно для каждой пары сигналов. Разрешение работы МТЗ НН2(СН) будет происходить по паре сигналов ОНМ с максимальным током. Углы максимальной чувствительности ОНМ в сторону трансформатора и в сторону сети задаются уставками \textbf{«Фмч ОНМ – В Т»} и \textbf{«Фмч ОНМ – В сеть»}. Угол максимальной чувствительности откладывается от линейного напряжения, за положительное направление принято вращение против часовой стрелки. Для ОНМ направленного в трансформатор, рекомендуется выбирать угол максимальной чувствительности в диапазоне 15º…60º. Для ОНМ направленного в сеть, рекомендуется выбирать угол максимальной чувствительности в диапазоне 195º…250º.  Величина тока работы ОНМ 0,05*Iном, величина напряжения работы ОНМ 0,01*Uном.

\medskip
\begin{figure}[H]
\centering
\includegraphics[width=1\textwidth,height=1\textheight,keepaspectratio]{\fbpath/06.01 ДЗТ/МТЗ НН2/img/ONM.pdf}
\captionof{figure}{Характеристики ОНМ МТЗ СН:\\ \textit{а} -- ОНМ включен на Ia и Ubc; \textit{б} -- ОНМ включен на Ic и Uab}
\label{pic:ONM}
\end{figure}

\item
Режим направленности для каждой ступени задается с помощью уставок \textbf{«Направ МТЗ-1 НН2(СН)»} и \textbf{«Направ МТЗ-2 НН2(СН)»}:
\begin{itemize}
\item \textit{«Ненапр»} - МТЗ СН без учета направления;
\item \textit{«В Т»} - МТЗ СН направлена в сторону трансформатора;
\item \textit{«В сеть»} - МТЗ СН  направлена в сторону сети.
\end{itemize}

\item
Режим работы направленных ступеней МТЗ НН2(СН) при неисправности ТН задается с помощью уставки \textbf{«КЦН ОНМ МТЗ»}:
\begin{itemize}
\item \textit{«Выведено»} - Режим работы ступени МТЗ НН2(СН) не зависит от БНН;
\item \textit{«РНМ»} - Ступень МТЗ НН2(СН) переводится в ненаправленный режим работы;
\item \textit{«МТЗ»} - Вывод направленных ступеней МТЗ НН2(СН).
\end{itemize}

\item
Алгоритм работы ОНМ СН приведен на рисунке \ref{pic:ONM_sch}. 

\medskip
\begin{figure}[H]
\centering
\includegraphics[width=1\textwidth,height=1\textheight,keepaspectratio]{\fbpath/06.01 ДЗТ/МТЗ НН2/img/onm_sch.pdf}
\captionof{figure}{Функциональная схема логики ОНМ СН}
\label{pic:ONM_sch}
\end{figure}


\end{enumerate}